\documentclass[12pt]{article}
\usepackage[utf8]{inputenc}
\usepackage[T1]{fontenc}
\usepackage{hyperref}
\title{Towei Language Resources for K-12 Curriculum Development}
\author{Compiled by Claude Sonnet 5 (claude-sonnet-5) for Andr\'as Kornai}
\date{}
\begin{document}
\maketitle

\section{Introduction}
Towei (ISO 639-3: ttn; Glottocode towe1240) is a Western Pauwasi language of Papua, Indonesia. It is spoken in Towe Hitam village, Towe District, Pegunungan Bintang Regency. Speaker counts vary by source: about 200 (a 1975 figure cited on Wikipedia), 120 (ASJP) and 115 (Global Recordings Network). Very little has been published on the language.

\section{Primary Educational Resources}
The searches found no primers, textbooks or school programmes for Towei.

\section{Linguistic and Reference Materials}
\begin{itemize}
\item Wikipedia overview of Towei, which cites a 2004 SIL survey report by Theresia Wambaliau: \url{https://en.wikipedia.org/wiki/Towei_language}
\item Wikipedia overview of the Pauwasi languages: \url{https://en.wikipedia.org/wiki/Pauwasi_languages}
\item ASJP word-list database entry: \url{https://asjp.clld.org/languages/TOWE}
\item Endangered Languages Project entry: \url{https://endangeredlanguages.com/lang/1487?hl=en}
\end{itemize}

\section{Religious and Community Resources}
The searches found no scripture or community organisation for Towei. The Global Recordings Network page for ttn is listed below. It states that it has no audio recordings in the language.
\begin{itemize}
\item \url{https://globalrecordings.net/en/language/ttn}
\end{itemize}

\section{Digital and Audio Resources}
No audio, video or corpus resources were found.

\section{Curriculum Development Context}
Towei has very few speakers and, as far as the searches show, is spoken mainly by older adults. No teaching materials were found. Any curriculum would have to start with documentation and community work, and I found nothing that could serve as a ready-made base.
\end{document}